\section{Cross ratios}

\label{sec:crossratios}

A \emph{generalized cross-ratio} is a function $[ \cdot, \cdot, \cdot, \cdot ]$ on $\partial S^{(4)}$ satisfying:

\begin{enumerate}
\item $[a,b,c,d]=[c,d,a,b]$ for any ccw tuple $(a,b,c,d)$.
\item If $a,b,c,d,e$ appear ccw, $[a,b,d,e]=[a,b,c,e]+[a,c,d,e]$
\item $\Gamma$-invariance.
\end{enumerate}

Given a generalized cross-ratio $[ \cdot, \cdot, \cdot, \cdot ]$, let the period function of a curve $c$ on $\Sigma$ represented by a group element $g$, be defined as
\[
\ell(g) \coloneqq [g^-, x , g \cdot x, g^+ ]
\]
From the properties of a generalized cross-ratio, it is easy to see that $\ell(g)=\ell(G)$. \DMGtodo[inline]{insert proof here? we have it.}

$\ell$ induces thus a function on curves that can be extended linearly to multi-curves. In this section we prove that this function satisfies the smoothing property.
We note that this can be inferred a posteriori from \cite{BIPP21:PositiveCrossratios}, where they construct a geodesic current for every generalized cross-ratio. It follows from this that $\ell$ must satisfy smoothing. However, the proof we give here uses only the definition of generalized cross-ratio and thus subsumes the work in \cite{BIPP21:PositiveCrossratios} under the scope of our theorem. Moreover, the proof also provides a framework on what cases to consider if one desires to check the smoothing property by working in the universal cover, as opposed to doing surgeries on the surface.
THis will also be helpful in establishing the generalized max-smoothing property for measured laminations.

We divide the proof into several cases, depending on the type of smoothing (connected vs disconnected) and the relative position of certain axes of the resolutions.

\subsection{Disconnected smoothing}

In this situation, we will see that even though there is only one case to consider, the proof is a bit more involved, relying on some tricks using the $\Gamma$ invariance property of the cross-ratio.

\begin{lemma}
Given two $C_1$ and $C_2$ represented by group elements $a$ and $b$ with an essential crossing, we prove that

\[
\ell(a) + \ell(b) \geq \ell(ab).
\]
\end{lemma}
\begin{proof}

The proof consists on decomposing $[(ba)^-, (Ab)^+, (ba) \cdot (Ab)^+, (ba)^+]$ into boxes $(1,I), (1,II), (2,I), (2,II), (3,I)$ and $(3, II)$ and showing, by splitting these boxes into sub-boxes and applying $b$ and $b^{-1}$ conveniently, that each one of these sub-boxes can be contained in either $[a^-, (Ab)^+, a \cdot (Ab)^+, a^+]$ or $[b^-, (ba)^+, b \cdot (bA)^+, b^+]$ without overlaps. We also note that $(ba)(Ab)(ba)^{-1}=b(bA)b^{-1}$, which will play a key role in the pictures below.
In the pictures, the bold font in the caption indicates that the result is final, and will not be perturbed further.

\begin{figure}[h!]
     \centering
     \begin{subfigure}[b]{0.45\textwidth}
\centering{
\resizebox{80mm}{!}{\svgfig{cr0}}
    \caption{Initial configuration}\label{fig:cr0}}
     \end{subfigure}
     \hfill
     \begin{subfigure}[b]{0.45\textwidth}
\centering{
\resizebox{80mm}{!}{\svgfig{cr1a1}}
    \caption{Box (1, I)}\label{fig:cr1a1}}
     \end{subfigure}
\end{figure}

\newpage

\begin{figure}[h!]
     \centering
     \begin{subfigure}[b]{0.45\textwidth}
\centering{
\resizebox{80mm}{!}{\svgfig{cr1a2}}
    \caption{\textbf{$b^{-1}$ applied to box $(1,I)$. Keep only one portion.}}\label{fig:cr1a2}}
     \end{subfigure}
     \hfill
     \begin{subfigure}[b]{0.45\textwidth}
\centering{
\resizebox{80mm}{!}{\svgfig{cr1a5}}
    \caption{Split from previous}\label{fig:cr1a5}}
     \end{subfigure}
\end{figure}

\begin{figure}[h!]
     \centering
     \begin{subfigure}[b]{0.45\textwidth}
\centering{
\resizebox{80mm}{!}{\svgfig{cr1a3}}
    \caption{\textbf{$b$ applied to previous}}\label{fig:cr1a3}}
     \end{subfigure}
     \hfill
     \begin{subfigure}[b]{0.45\textwidth}
\centering{
\resizebox{80mm}{!}{\svgfig{cr1a4}}
    \caption{$b$ applied to $(1,II)$}\label{fig:cr1a4}}
     \end{subfigure}
\end{figure}

\begin{figure}[h!]
     \centering
     \begin{subfigure}[b]{0.45\textwidth}
\centering{
\resizebox{80mm}{!}{\svgfig{cr1a6}}
    \caption{\textbf{Split one portion}}\label{fig:cr1a6}}
     \end{subfigure}
\hfill
     \begin{subfigure}[b]{0.45\textwidth}
\centering{
\resizebox{80mm}{!}{\svgfig{cr1a5}}
    \caption{\textbf{Split other portion, apply $b^{-1}$}}\label{fig:cr1a5}}
     \end{subfigure}
\end{figure}


\begin{figure}[h!]
     \centering
     \begin{subfigure}[b]{0.45\textwidth}
\centering{
\resizebox{80mm}{!}{\svgfig{cr2a1}}
    \caption{\textbf{$(2,I)$}}\label{fig:cr2a1}}
     \end{subfigure}
     \hfill
     \begin{subfigure}[b]{0.45\textwidth}
\centering{
\resizebox{80mm}{!}{\svgfig{cr2b1}}
    \caption{$(2,II)$}\label{fig:cr2b1}}
     \end{subfigure}
\end{figure}

\begin{figure}[h!]
     \centering
     \begin{subfigure}[b]{0.45\textwidth}
\centering{
\resizebox{80mm}{!}{\svgfig{cr2b2}}
    \caption{Split portion of $(2,II)$}\label{fig:cr2b2}}
     \end{subfigure}
     \hfill
     \begin{subfigure}[b]{0.45\textwidth}
\centering{
\resizebox{80mm}{!}{\svgfig{cr2b3}}
    \caption{\textbf{Apply $b$ to previous}}\label{fig:cr2b3}}
     \end{subfigure}
\end{figure}

\begin{figure}[h!]
     \centering
     \begin{subfigure}[b]{0.45\textwidth}
\centering{
\resizebox{80mm}{!}{\svgfig{cr3a1}}
    \caption{Box $(3,I)$}\label{fig:cr3a1}}
     \end{subfigure}
     \hfill
     \begin{subfigure}[b]{0.45\textwidth}
\centering{
\resizebox{80mm}{!}{\svgfig{cr3a2}}
    \caption{ $b^{-1}$ applied to previous}\label{fig:cr3a2}}
     \end{subfigure}
\end{figure}

\begin{figure}[h!]
     \centering
     \begin{subfigure}[b]{0.45\textwidth}
\centering{
\resizebox{80mm}{!}{\svgfig{cr3b1}}
    \caption{Box $(3,II)$}\label{fig:cr3b1}}
     \end{subfigure}
     \hfill
     \begin{subfigure}[b]{0.45\textwidth}
\centering{
\resizebox{80mm}{!}{\svgfig{cr3b2}}
    \caption{Split previous}\label{fig:cr3b2}}
     \end{subfigure}
\end{figure}

\begin{figure}[h!]
     \centering
     \begin{subfigure}[b]{0.45\textwidth}
\centering{
\resizebox{80mm}{!}{\svgfig{cr3b3}}
    \caption{\textbf{$b$ applied to previous}}\label{fig:cr3b3}}
     \end{subfigure}
     \hfill
     \begin{subfigure}[b]{0.45\textwidth}
\centering{
\resizebox{80mm}{!}{\svgfig{cr3b4}}
    \caption{\textbf{Remaining portion of $(3,II)$}}\label{fig:cr3b4}}
     \end{subfigure}
\end{figure}

\end{proof}

\subsection{Connected smoothing}

If $C$ represents a curve with a self-intersection decomposing as $ab$, we have several cases to consider, depending on the relative position of $a$ and $b$, and on the position of $aB$ and $Ba$ (intersecting or non-intersecting).

\subsubsection{Parallel $a$ and $b$}

In the case when $a$ and $b$ are parallel, there are two cases to consider.

\begin{figure}[h!]
     \centering
\resizebox{80mm}{!}{\svgfig{c1}}
    \caption{$a$ and $b$ parallel, case 1 ($aB$ and $Ba$ do not intersect)}\label{fig:c1}
\end{figure}


\begin{figure}[h!]
     \centering
\resizebox{80mm}{!}{\svgfig{c2}}
    \caption{$a$ and $b$ parallel, case 2 ($aB$ and $Ba$  intersect)}\label{fig:c2}
\end{figure}

\begin{lemma}
We have
\[
\ell(ab) \geq \ell(aB)
\]
and
\[
\ell(ab) \geq \ell(a) + \ell(b).
\]
\end{lemma}
\begin{proof}
We start by analyzing case 1 first.
For the first inequality compare Figure \ref{fig:c1}. The proof follows by noticing that $AB \cdot (baB)^- = a^-$,
    and
    \[
    [baB^+, ba^+, ba^-, a-] = [ba^-, a^-, baB^-, ba^+]=\ell(ba).
    \]
    Since $(aB^+, aB^-) \subset (AB^+, AB^-)$, we have
    \[
    \ell(aB)= [aB^-, (baB)^-, a^-, aB^+] \leq [ba^-, a^-, baB^-, ba^+]=\ell(ba).
    \]
For the second inequality, we can write, using the flip relation and the additive relation of the cross-ratio
\begin{align*}
\ell(ba)&=[ba^-, a^-, b \cdot a^-, ba^+] = [b \cdot a^-, ba^+,ba^-, a^-] =\\
&[b \cdot a^-, ba^+,b^+, a^-] +  [b \cdot a^-, b^+,b^-, a^-] + \\
&[b \cdot a^-, b^-,ab^-, a^-] + [b \cdot a^-, ab^-,ba^-, a^-].
\end{align*}
Therefore, by using the addivity property repeatedly, we get
\begin{align*}
\ell(ba) &\geq  [b \cdot a^-, b^+,b^-, a^-] + [b \cdot a^-, ab^-,ba^-, a^-] =\\
&[b \cdot a^-, b^+,b^-, a^-] +[ba^-, a^-, b \cdot a^-, ab^-]\\
&[b \cdot a^-, b^+,b^-, a^-] +[ba^-, a^-, a^+, ab^-] +  [ba^-, a^+, b \cdot a^-, ab^-] \geq \\
& [b \cdot a^-, b^+,b^-, a^-] +[ba^-, a^-, a^+, ab^-] = \\
& [b \cdot a^-, b^+,b^-, a^-] +[ba^-, A^+, A^-, ab^-] = \\
& \ell(b) + \ell(A) = \ell(b) + \ell(a)
\end{align*}
where in the last equation we used the fact that $\ell(a)=\ell(A)$.
The same proof follows verbatim in case 2 for the second inequality, since we have not used the position of the axes $aB$ and $Ba$.

It remains to prove the first inequality for case 2.
\end{proof}



%%% Local Variables:
%%% mode: latex
%%% TeX-master: "Intersections"
%%% End: